\section{Pipeline de données}

Les données affichées par l'application proviennent de fichiers JSON statiques, générés par des scripts Python à partir de sources externes (dictionnaire Goroh en ligne) et de phrases annotées manuellement.

\subsection{Vue d'ensemble}

\begin{lstlisting}[title=Pipeline de generation des donnees]
phrases_a_traiter.json    (phrases annotees avec spans .ukr)
        |
        v
build_entries_2_(nooj_opt).py   (scraping + parsing morphologique)
        |
        v
    out.json              (nouvelles entrees generees)
        |
        v  (fusion manuelle)
  static/data.json        (dictionnaire complet, utilise par l'app)
  static/phrases.json     (phrases d'exemple)
\end{lstlisting}

\subsection{Fichiers du pipeline}

Tous situés dans \texttt{outil\_python/}~:

\begin{description}
  \item[\texttt{build\_entries\_2\_(nooj\_opt).py}] Script principal. Pour chaque mot ukrainien trouvé dans les phrases à traiter~:
  \begin{enumerate}
    \item Vérifie s'il existe déjà dans \texttt{data.json} (dédoublement).
    \item Cherche le lemme et la catégorie grammaticale.
    \item Scrape le dictionnaire Goroh (\texttt{goroh.ua}) pour obtenir les tables de déclinaison/conjugaison.
    \item Parse les tables HTML en structure JSON normalisée.
    \item Génère le champ \texttt{base\_html} (span canonique avec \texttt{data-info}).
    \item Écrit le résultat dans \texttt{out.json}.
  \end{enumerate}

  \item[\texttt{phrases\_a\_traiter.json}] Fichier d'entrée : liste de phrases avec leurs spans \texttt{.ukr} annotés (traduction, références, remarques).

  \item[\texttt{out.json}] Fichier de sortie : nouvelles entrées prêtes à être fusionnées dans \texttt{static/data.json}.

  \item[\texttt{entries\_report.html}] Rapport HTML des entrées traitées (pour vérification manuelle).
\end{description}

\subsection{Format des données générées}

Chaque entrée dans \texttt{out.json} suit le format attendu par l'application (voir annexe~\ref{sec:data-schema}). Points importants~:

\begin{itemize}
  \item \textbf{Ordre des champs} : \texttt{cas} (ou \texttt{inf}/\texttt{conj} pour les verbes), puis \texttt{genre} (noms), \texttt{asp} (verbes), \texttt{nooj} (optionnel), \texttt{base\_html}, \texttt{phrases}.
  \item \textbf{Collisions de lemmes} : si deux mots ont le même lemme mais des catégories différentes, le suffixe \texttt{\_\_pos} est ajouté (ex: \texttt{один\_\_card} vs \texttt{один\_\_adj}).
  \item \textbf{Format canonique de \texttt{base\_html}} :
\end{itemize}

\begin{lstlisting}[language=html]
<span class="ukr" data-info="LEMME;CATEGORIE;...">TEXTE</span>
\end{lstlisting}

\subsection{Mise à jour des données}

Le processus de mise à jour est semi-automatique~:

\begin{enumerate}
  \item Ajouter les nouvelles phrases dans \texttt{phrases\_a\_traiter.json}.
  \item Exécuter le script Python : \texttt{python build\_entries\_2\_(nooj\_opt).py}.
  \item Vérifier \texttt{entries\_report.html} pour les erreurs.
  \item Fusionner manuellement \texttt{out.json} dans \texttt{static/data.json}.
  \item Ajouter les phrases dans \texttt{static/phrases.json}.
  \item Rebuilder l'application : \texttt{npm run build}.
\end{enumerate}
